% Paquets généraux
\documentclass[a4paper,12pt,titlepage]{article}
\usepackage[T1]{fontenc}
\usepackage[utf8]{inputenc}
\usepackage[french]{babel}
\usepackage[gen]{eurosym}
%\usepackage[dvips]{graphicx}
\usepackage{fancyhdr}
\usepackage{pdfpages} 
\usepackage{multido}
\usepackage{moreverb}
\usepackage{hyperref}
%\usepackage{textcomp}
\usepackage{verbatim}
\usepackage{moreverb}
\usepackage{listings}
\usepackage{minted}
\usepackage{eso-pic}
\usepackage{enumitem}
\usepackage{comment}
\usepackage{boxedminipage}
\usepackage[french,onelanguage, boxruled,linesnumbered]{algorithm2e}


\newcommand{\auteurun}{Juliette Genzmer}
\newcommand{\auteurdeux}{Willie Robert}
\newcommand{\auteurtrois}{Renaud Costadoat}
\newcommand{\institute}{Lycée Dorian}
\newtheorem{solution}{Solution}


\newcommand{\nom}{Initiation à l'algorithmie}
\newcommand{\sequence}{1}
\newcommand{\numero}{01}
\newcommand{\type}{DS}
\newcommand{\descrip}{Initiation à l'algorithmie}


\usepackage{color}
\usepackage{xcolor}
\usepackage{colortbl}
\usepackage{helvet}
\renewcommand{\familydefault}{\sfdefault}
\usepackage{amsfonts}
\usepackage{amsmath}
%\usepackage{xspace}
\usepackage{varioref}
\usepackage{tabularx}
%\usepackage{floatflt}
\usepackage{graphics}
\usepackage{wrapfig}
\usepackage{textcomp}
\usepackage{tikz}
\usepackage{wrapfig}
\usepackage{gensymb}
\usepackage{ifthen}
\usepackage{cancel}
\usepackage{etoolbox}
\usepackage{multirow}
%\usepackage{boxedminipage}
\definecolor{gris25}{gray}{0.75}
\definecolor{bleu}{RGB}{18,33,98}
\definecolor{bleuf}{RGB}{42,94,171}
\definecolor{bleuc}{RGB}{231,239,247}
\definecolor{rougef}{RGB}{185,18,27}
\definecolor{rougec}{RGB}{255,188,204}%255,230,231
\definecolor{vertf}{RGB}{103,126,82}
\definecolor{vertc}{RGB}{220,255,191}
\definecolor{forestgreen}{rgb}{0.13,0.54,0.13}
\definecolor{blcr}{rgb}{0.59,0.69,0.84}
\definecolor{blfr}{rgb}{0.32,0.51,0.75}
\definecolor{orfr}{rgb}{0.90,0.42,0.15}
\definecolor{orcr}{rgb}{0.90,0.65,0.50}
\definecolor{orangef}{rgb}{0.659,0.269,0.072}
\definecolor{orange}{rgb}{0.58,0.35,0.063}
\definecolor{orangec}{rgb}{0.43,0.32,0.25}
\definecolor{rcorrect}{rgb}{0.6,0,0}
\definecolor{sequence}{rgb}{0.75,0.75,0.75}
\definecolor{competences}{rgb}{0.61,0.73,0.35}
\definecolor{grisf}{HTML}{222222}
\definecolor{grisc}{HTML}{636363}
\definecolor{normal}{HTML}{4087c4}
\definecolor{info}{HTML}{5bc0de}
\definecolor{success}{RGB}{92,184,92}
\definecolor{warning}{RGB}{240,173,78}
\definecolor{danger}{RGB}{217,83,79}
\hypersetup{                    % parametrage des hyperliens
    colorlinks=true,                % colorise les liens
    breaklinks=true,                % permet les retours à la ligne pour les liens trop longs
    urlcolor= blfr,                 % couleur des hyperliens
    linkcolor= orange,                % couleur des liens internes aux documents (index, figures, tableaux, equations,...)
    citecolor= forestgreen                % couleur des liens vers les references bibliographiques
    }

% Mise en page
\pagestyle{fancy}

\setlength{\hoffset}{-18pt}

\setlength{\oddsidemargin}{0pt} 	% Marge gauche sur pages impaires
\setlength{\evensidemargin}{0pt} 	% Marge gauche sur pages paires
\setlength{\marginparwidth}{00pt} 	% Largeur de note dans la marge
\setlength{\headwidth}{481pt} 	 	% Largeur de la zone de tête (17cm)
\setlength{\textwidth}{481pt} 	 	% Largeur de la zone de texte (17cm)
\setlength{\voffset}{-18pt} 		% Bon pour DOS
\setlength{\marginparsep}{7pt}	 	% Séparation de la marge
\setlength{\topmargin}{-30pt} 		% Pas de marge en haut
\setlength{\headheight}{55pt} 		% Haut de page
\setlength{\headsep}{20pt} 		% Entre le haut de page et le texte
\setlength{\footskip}{30pt} 		% Bas de page + séparation
\setlength{\textheight}{700pt} 		% Hauteur de l'icone zone de texte (25cm)
\setlength\fboxrule{1 pt}
\renewcommand{\baselinestretch}{1}
\setcounter{tocdepth}{1}
\newcommand{\cadre}[2]
{\fbox{
  \begin{minipage}{#1\linewidth}
   \begin{center}
    #2\\
   \end{center}
  \end{minipage}
 }
}

\newcounter{num_quest} \setcounter{num_quest}{0}
\newcounter{num_rep} \setcounter{num_rep}{0}
\newcounter{num_cor} \setcounter{num_cor}{0}

\newcommand{\question}[1]{\refstepcounter{num_quest}\par
~\ \\ \parbox[t][][t]{0.15\linewidth}{\textbf{Question \arabic{num_quest}}}\parbox[t][][t]{0.85\linewidth}{#1\label{q\the\value{num_quest}}}\par\par
}



\newcommand{\reponse}[0]{\refstepcounter{num_rep}\par
~\ \\ \parbox[t][][t]{0.15\linewidth}{\textbf{Question \arabic{num_rep}}}}

\newcommand{\cor}
{\refstepcounter{num_cor}
\noindent
\rule{\linewidth}{.5pt}
\textbf{Question \arabic{num_cor}:} \\
}



% En tête et pied de page
\lhead{\nom}
\rhead{\includegraphics[width=2cm]{../../../img/logo}}
\lfoot{David Aubert, Renaud Costadoat}
\rfoot{Page \thepage}
\cfoot{}

\newlength{\RoundedBoxWidth}
\newsavebox{\GrayRoundedBox}
\newenvironment{GrayBox}[1][\dimexpr\textwidth-4.5ex]%
   {\setlength{\RoundedBoxWidth}{\dimexpr#1}
    \begin{lrbox}{\GrayRoundedBox}
       \begin{minipage}{\RoundedBoxWidth}}%
   {   \end{minipage}
    \end{lrbox}
    \begin{center}
    \begin{tikzpicture}%
       \draw node[draw=bleuf,fill=bleuc,rounded corners,%
             inner sep=2ex,text width=\RoundedBoxWidth]%
             {\usebox{\GrayRoundedBox}};
    \end{tikzpicture}
    \end{center}}

\fancypagestyle{correction}{%
  \fancyhf{}
  \lhead{\colorbox{danger}{\begin{minipage}{0.65\paperwidth} \textcolor{white}{\textbf{Correction}} \end{minipage}} }
  \rhead{\includegraphics[width=2cm]{../../../img/logo}}
  \lfoot{Juliette Genzmer, Willie Robert, Renaud Costadoat}
  \rfoot{\colorbox{danger}{\begin{minipage}{0.3\paperwidth} \begin{flushright}\textcolor{white}{\textbf{Correction}}\end{flushright} \end{minipage}} }}

\renewcommand{\footrulewidth}{0.4pt}


\newcommand{\BackgroundPic}{%
\put(0,0){%
\parbox[b][\paperheight]{\paperwidth}{%
\vfill
\begin{center}
\hspace{0.5cm}\vspace{0.5cm}
\includegraphics[width=\paperwidth,height=\paperheight,%
keepaspectratio]{../../../img/fond5}%
\end{center}
\vfill
}}}

\newcommand{\goforum}{
\begin{figure}[ht!]
\begin{center}
 \includegraphics[width=0.7\linewidth]{../../../img/go_forum}
\end{center}
\label{go_forum}
\caption{J'pète les plombs}
\end{figure}}

\newcommand{\BackgroundPicdeux}{%
\put(25,-30){%
\parbox[b][\paperheight]{\paperwidth}{%
\vfill
\begin{center}
\includegraphics[width=\paperwidth,height=\paperheight,%
keepaspectratio]{../../../img/fond6}%
\end{center}
\vfill
}}}

\setenumerate[1]{align=left,label=\arabic*}
\setenumerate[2]{before=\stepcounter{enumi},label*=.\arabic*,leftmargin=1.2em,align=left}

\begin{document}


\pagestyle{fancy}

\AddToShipoutPicture{\BackgroundPicdeux}

\begin{center}
{\Large\bf {\type} \no {\numero} -- \descrip}
\end{center}

\SetKw{KwFrom}{de}

\begin{boxedminipage}{.9\textwidth}
\begin{itemize}
 \item Faire tous les exercices dans un fichier {NomPrenom.py} à sauvegarder,
 \item mettre en commentaire l'exercice et la question traités (ex: \# Exercice 1),
 \item ne pas oublier pas de commenter ce qui est fait dans votre code (ex: \# Je crée une fonction pour calculer la racine d'un nombre),
 \item il faut vérifier avant de partir que le code peut s'exécuter et qu'il affiche les résultats que vous attendez. Les lignes de code qui doivent s'exécuter sont décommentées.
\end{itemize}
\end{boxedminipage}

%\title{Devoir d'informatique en Python \\ \large Simulation d'une billetterie de concert}

\section*{Contexte}
La mise en vente des billets pour un concert très attendu approche. Les années précédentes, pour des tournées similaires, le site de billetterie officiel s'est effondré en quelques minutes à cause du nombre de connexions simultanées. Pour éviter cela, on souhaite programmer une version simplifiée d'un système de vente de billets en ligne.

\begin{figure}[!ht]
	\begin{center}
		\includegraphics[width=0.4\linewidth]{img/concert.jpg}
	\end{center}
\end{figure}

La billetterie propose au départ~:
\begin{itemize}
    \item Fosse : 133.50 € (stock initial : 5),
    \item Gradin : 232.5 € (stock initial : 5),
    \item Carré Or : 298.5 € (stock initial : 5).
\end{itemize}

\bigskip
L'objectif est de réaliser un programme en Python qui simule les achats d'un spectateur.

\section{Choix de la catégorie}

\begin{minted}[linenos]{python}
import random

# Stocks de départ
fosse = 5
gradin = 5
carre_or = 5

0 = achats

print("Bienvenue à la billetterie en ligne !)

# Boucle tant qu'il reste des places
while fosse > 0 or gradin > 0 or carre_or > 0
    print("\n=== Catégories disponibles ===")
    print("1 - Fosse (133.50 euros) [stock :", fosse, "]")
    print("2 - Gradin (232.5 euros) [stock :", gradin, "]")
    print("3 - Carré Or (298.5 euros) [stock :", carre_or, "]")
    print("0 - Quitter")

    choix = int(input("Sélection ? "))

    if choix = 0:
        print("Merci, à bientôt !")
        break
    else:
        if choix == 1:  # Cas Fosse
            print("Place en fosse demandée")

print("\n=== Fin du programme ===")
print("Nombre total de billets vendus :", achats)
\end{minted}

Ce programme permet de choisir la catégorie de place en entrant un entier au clavier. La suite du programme n'étant pas codée, il faut entrer 0 pour quitter et retrouver la console.

\question{\textbf{Recopier} le code, \textbf{corriger} les 4 erreurs qui se sont glissées à l'intérieur (penser à lire les erreurs sur la console) et l'\textbf{exécuter} afin de vérifier que la saisie du choix est possible.}

\newpage

\section{Paiement et confirmation}

Nous nous intéressons maintenant au cas de figure où le spectateur a entré 1 avec son clavier afin de sélectionner une place en fosse.

Le programme doit maintenant :
\begin{enumerate}
    \item Vérifier que le stock est suffisant. Si la catégorie est épuisée, afficher un message d'erreur.
    \item Demander combien d'argent le spectateur insère.
    \item Vérifier que la somme est suffisante :
    \begin{itemize}
        \item si oui, calculer la monnaie à rendre,
        \item si non, refuser la vente et rendre l'argent.
    \end{itemize}
    \item Modifier la quantité de places en fosse en stock,
    \item Modifier le nombre total de billets vendus.
\end{enumerate}

\begin{minted}[linenos]{python}
if choix == 1:
    print("Place en fosse demandée")
    if .........:
        print("Plus de places en fosse.")
    else:
        # Argent inséré
        argent = int(input("Montant inséré : "))
        if argent ......:
            print("Montant insuffisant.")
        else:
            print("Vous avez acheté une place en fosse.")
            monnaie = ..........
            if monnaie > 0:
                print("Monnaie rendue :", monnaie, "euros")
            fosse .....
            achats .....
\end{minted}

\question{\textbf{Recopier} le code à sa bonne place dans le script (attention à l'indentation), le \textbf{compléter} et l'\textbf{exécuter} afin de vérifier le fonctionnement lors du choix d'une place en fosse.}

\question{A partir de cet extrait, \textbf{coder} la suite du script permettant la gestion de l'achat d'une place en gradin ou en carré or. Il faudra aussi prévoir un message d'erreur si le spectateur effectue une sélection erronée.}

\newpage

\section{Organisation en fonctions}

Afin d'améliorer la lisibilité du code, il est préférable d'utiliser les deux fonctions :
\begin{itemize}
    \item \texttt{affichage\_menu(fosse, gradin, carre\_or)},
    \item \texttt{vendre\_billet(nom, stock, prix, achats)}.
\end{itemize}

La suite du code pourrait alors s'écrire comme suit.
\begin{minted}[linenos]{python}
    if choix == 0:
        print("Merci, à bientôt !")
        break
    else:
        if choix == 1:  # Cas Fosse
            fosse,achats=vendre_billet('fosse',fosse,133.50,achats)
        elif choix == 2:  # Cas Gradin
            gradin,achats=vendre_billet('gradin',gradin,232.5,achats)
        elif choix == 3:  # Cas Carré Or
            carre_or,achats=vendre_billet('carre or',carre_or,298.5,achats)
        else:
            print("Choix invalide.")
\end{minted}

\question{Écrire les deux fonctions \texttt{affichage\_menu} et \texttt{vendre\_billet} et les tester à l'aide du code fourni.}

\section{Offre commerciale}

\question{\textbf{Ajouter} une option \og surclassement offert\fg, à la fonction \texttt{vendre\_billet}, qui une fois sur 20 surclasse gratuitement le spectateur en Carré Or (utiliser l'opération \texttt{random.randrange(20) == 0}).}

\begin{center}
\Large{FIN}
\end{center}

\cleardoublepage

\ifdef{\public}{\end{document}}{\pagestyle{correction}}

\begin{center}
\Large{Correction}
\end{center}

\reponse{}

\begin{minted}{python}
import random

# Stocks de départ
fosse = 5
gradin = 5
carre_or = 5

achats = 0

print("Bienvenue à la billetterie en ligne !")

# Boucle tant qu'il reste des places
while fosse > 0 or gradin > 0 or carre_or > 0:
    print("\n=== Catégories disponibles ===")
    print("1 - Fosse (133.50 euros) [stock :", fosse, "]")
    print("2 - Gradin (232.5 euros) [stock :", gradin, "]")
    print("3 - Carré Or (298.5 euros) [stock :", carre_or, "]")
    print("0 - Quitter")

    choix = int(input("Sélection ? "))

    if choix == 0:
        print("Merci, à bientôt !")
        break
    else:
        # Cas Fosse
        if choix == 1:
            print("Place en fosse demandée")

print("\n=== Fin du programme ===")
print("Nombre total de billets vendus :", achats)
\end{minted}

\reponse{}

\begin{minted}{python}
import random

# Stocks de départ
fosse = 5
gradin = 5
carre_or = 5

achats = 0

print("Bienvenue à la billetterie en ligne !")

# Boucle tant qu'il reste des places
while fosse > 0 or gradin > 0 or carre_or > 0:
    print("\n=== Catégories disponibles ===")
    print("1 - Fosse (133.50 euros) [stock :", fosse, "]")
    print("2 - Gradin (232.5 euros) [stock :", gradin, "]")
    print("3 - Carré Or (298.5 euros) [stock :", carre_or, "]")
    print("0 - Quitter")

    choix = int(input("Sélection ? "))

    if choix == 0:
        print("Merci, à bientôt !")
        break
    else:
        # Cas Fosse
        if choix == 1:
            print("Place en fosse demandée")
            if fosse == 0:
                print("Plus de places en fosse.")
            else:
                argent = int(input("Montant inséré : "))
                if argent < 133.50:
                    print("Montant insuffisant.")
                else:
                    print("Vous avez acheté une place en fosse.")
                    monnaie = argent - 133.50
                    if monnaie > 0:
                        print("Monnaie rendue :", monnaie, "euros")
                    fosse -= 1
                    achats += 1

print("\n=== Fin du programme ===")
print("Nombre total de billets vendus :", achats)
\end{minted}

\reponse{}

\begin{minted}{python}
import random

# Stocks de départ
fosse = 5
gradin = 5
carre_or = 5

achats = 0

print("Bienvenue à la billetterie en ligne !")

# Boucle tant qu'il reste des places
while fosse > 0 or gradin > 0 or carre_or > 0:
    print("\n=== Catégories disponibles ===")
    print("1 - Fosse (133.50 euros) [stock :", fosse, "]")
    print("2 - Gradin (232.5 euros) [stock :", gradin, "]")
    print("3 - Carré Or (298.5 euros) [stock :", carre_or, "]")
    print("0 - Quitter")

    choix = int(input("Sélection ? "))

    if choix == 0:
        print("Merci, à bientôt !")
        break
    else:
        # Cas Fosse
        if choix == 1:
            print("Place en fosse demandée")
            if fosse == 0:
                print("Plus de places en fosse.")
            else:
                argent = int(input("Montant inséré : "))
                if argent < 133.50:
                    print("Montant insuffisant.")
                else:
                    print("Vous avez acheté une place en fosse.")
                    monnaie = argent - 133.50
                    if monnaie > 0:
                        print("Monnaie rendue :", monnaie, "euros")
                    fosse -= 1
                    achats += 1

        # Cas Gradin
        elif choix == 2:
            if gradin == 0:
                print("Plus de places en gradin.")
            else:
                argent = int(input("Montant inséré : "))
                if argent < 232.5:
                    print("Montant insuffisant.")
                else:
                    print("Vous avez acheté une place en gradin.")
                    monnaie = argent - 232.5
                    if monnaie > 0:
                        print("Monnaie rendue :", monnaie, "euros")
                    gradin -= 1
                    achats += 1

        # Cas Carré Or
        elif choix == 3:
            if carre_or == 0:
                print("Plus de places en carré or.")
            else:
                argent = int(input("Montant inséré : "))
                if argent < 298.5:
                    print("Montant insuffisant.")
                else:
                    print("Vous avez acheté une place en carré or.")
                    monnaie = argent - 298.5
                    if monnaie > 0:
                        print("Monnaie rendue :", monnaie, "euros")
                    carre_or -= 1
                    achats += 1
        else:
            print("Choix invalide.")

print("\n=== Fin du programme ===")
print("Nombre total de billets vendus :", achats)
\end{minted}

\reponse{}

\begin{minted}{python}
import random

def affichage_menu(fosse, gradin, carre_or):
    print("\n=== Catégories disponibles ===")
    print("1 - Fosse (133.50 euros) [stock :", fosse, "]")
    print("2 - Gradin (232.5 euros) [stock :", gradin, "]")
    print("3 - Carré Or (298.5 euros) [stock :", carre_or, "]")
    print("0 - Quitter")

def vendre_billet(nom, stock, prix, achats):

    if stock == 0:
        print("Plus de places en "+nom+".")
    else:
        argent = int(input("Montant inséré : "))
        if argent < prix:
            print("Montant insuffisant.")
        else:
            print("Vous avez acheté une place en "+nom+".")
            monnaie = argent - prix
            if monnaie > 0:
                print("Monnaie rendue :", monnaie, "euros")
            stock -= 1
            achats += 1
    return stock, achats

# Stocks de départ
fosse = 5
gradin = 5
carre_or = 5

achats = 0

print("Bienvenue à la billetterie en ligne !")

# Boucle tant qu'il reste des places
while fosse > 0 or gradin > 0 or carre_or > 0:
    affichage_menu(fosse, gradin, carre_or)

    choix = int(input("Sélection ? "))

    if choix == 0:
        print("Merci, à bientôt !")
        break
    else:
        # Cas Fosse
        if choix == 1:
            fosse,achats=vendre_billet('fosse',fosse,133.50,achats)
        # Cas Gradin
        elif choix == 2:
            gradin,achats=vendre_billet('gradin',gradin,232.5,achats)
        # Cas Carré Or
        elif choix == 3:
            carre_or,achats=vendre_billet('carre or',carre_or,298.5,achats)
        else:
            print("Choix invalide.")

print("\n=== Fin du programme ===")
print("Nombre total de billets vendus :", achats)
\end{minted}

\reponse{}

\begin{minted}{python}
import random

def affichage_menu(fosse, gradin, carre_or):
    print("\n=== Catégories disponibles ===")
    print("1 - Fosse (133.50 euros) [stock :", fosse, "]")
    print("2 - Gradin (232.5 euros) [stock :", gradin, "]")
    print("3 - Carré Or (298.5 euros) [stock :", carre_or, "]")
    print("0 - Quitter")

def vendre_billet(nom, stock, prix, achats):

    if stock == 0:
        print("Plus de places en "+nom+".")
    elif random.randrange(20) == 0:
        print("Coup de chance ! Vous êtes surclassé(e) en Carré Or, offert !")
        stock -= 1
        achats += 1
    else:
        argent = int(input("Montant inséré : "))
        if argent < prix:
            print("Montant insuffisant.")
        else:
            print("Vous avez acheté une place en "+nom+".")
            monnaie = argent - prix
            if monnaie > 0:
                print("Monnaie rendue :", monnaie, "euros")
            stock -= 1
            achats += 1
    return stock, achats

# Stocks de départ
fosse = 5
gradin = 5
carre_or = 5

achats = 0

print("Bienvenue à la billetterie en ligne !")

# Boucle tant qu'il reste des places
while fosse > 0 or gradin > 0 or carre_or > 0:
    affichage_menu(fosse, gradin, carre_or)

    choix = int(input("Sélection ? "))

    if choix == 0:
        print("Merci, à bientôt !")
        break
    else:
        # Cas Fosse
        if choix == 1:
            fosse,achats=vendre_billet('fosse',fosse,133.50,achats)
        # Cas Gradin
        elif choix == 2:
            gradin,achats=vendre_billet('gradin',gradin,232.5,achats)
        # Cas Carré Or
        elif choix == 3:
            carre_or,achats=vendre_billet('carre or',carre_or,298.5,achats)
        else:
            print("Choix invalide.")

print("\n=== Fin du programme ===")
print("Nombre total de billets vendus :", achats)
\end{minted}

\end{document}
